\documentclass[a4paper,10pt]{report}

\def\packagepath{../../preambule}   % path du package princial
\usepackage{\packagepath/preambule}    % utilisation du fichier de configuration

\def\level{Première }              % Classe
\def\course{NSI}              % Matière
\def\eval{Codage des entiers relatifs}

\def\visibleornot{visible}    % visible or invisible
\def\documentpath{./Sources_Latex}
\def\date{Mardi 08 Avril 2025}
\renewcommand{\arraystretch}{1}  % Ecart dans les tableau

\def\ptsexoA{3}
\def\ptsexoB{3}
\def\ptsexoC{2}
\def\ptsexoD{2}
\def\ptstotal{\ptsexoA+\ptsexoB}

\usepackage{hyperref}

\usetikzlibrary{shapes, arrows, positioning}
\usetikzlibrary{shapes.geometric, arrows, positioning, decorations.pathreplacing}

\tikzstyle{etat} = [draw, rounded corners=15pt, minimum width=2.5cm, minimum height=1.2cm, text centered, font=\bfseries]
\tikzstyle{nouveau} = [etat, fill=blue!40]
\tikzstyle{pret} = [etat, fill=green!60]
\tikzstyle{elu} = [etat, fill=yellow!60]
\tikzstyle{bloque} = [etat, fill=red!60]
\tikzstyle{termine} = [etat, fill=white]
\tikzstyle{fleche} = [thick, ->, >=stealth]

\usepackage{float} % Permet d'utiliser [H]
\usepackage[T1]{fontenc}
\begin{document}

%%%%%%%%%%%%%%%%%%%%%%%%%%%%%%%%%%%%%%%%%%%%%%%%%%ù

%\PageGardeSujetBac[Matiere = Numérique et Sciences Informatiques,
%                  Session = 2025,
%                  AffJour=false,
%                  Duree = {3 heures 30},
%                  DernierePage = \pageref{LastPage},
%                  ModeExamen=false
%                  ]
\renewcommand{\labelitemi}{\textbullet} %pour éviter les tirets dans les "itemize" qui apportent confusion avec le signe moins.
%% Début page de garde style BAC

   %%%%%%%%%%%%%%%%%%%%%%%%%%%%%%%%%%%%%%%%%%%%%%%%%%%%%%%%
%\PageGardeSujetBac[clés]

\pagestyle{DS_FP}          % Feuille de style fancy
\NomPrenomNote{}


\exods{\ptsexoA}

Répondre aux questions suivantes. Aucune justification n'est demandée.

\begin{enumerate}
    \item Avec 8 bits, combien de valeurs différentes peut-on coder? $2^8=256$ valeurs différentes.
    \item Quelles sont les valeurs entières naturelles minimale et maximale codable sur 8 bits? $0$ et $255$.
    \item Quelles sont les valeurs entières relatives minimale et maximale codable sur 8 bits? $-128$ et $127$.
\end{enumerate}


\exods{\ptsexoB}

\begin{enumerate}
    \item Ecrire l'algorithme de conversion d'un entier naturel compris entre 0 et 255 en binaire sur 8 bits. Vous pouvez écrire cet algorithme en pseudo-code ou en Python.
    
    \begin{answer}{1}{5.5 }
        \begin{verbatim}
def conversion_dec_to_bin(val_dec):
    ''' Convertit un entier naturel en binaire sur 8 bits.
    val_dec : int : entier naturel compris entre 0 et 255   
    return : str : représentation binaire sur 8 bits de l'entier naturel'''
    binaire = ''
    # On effectue la conversion de l'entier décimal en binaire
    for i in range(8):
        reste = val_dec % 2
        binaire = str(reste) + binaire
        val_dec = val_dec // 2
    
    return binaire
        
        \end{verbatim}
        
        
        
        
        \end{answer}

    \item Donner l'algorithme de conversion d'un entier relatif en binaire sur 8 bits.
    
    Vous citerez les différentes étapes de l'algorithme sans le traduire en Python.
    \begin{answer}{1}{10}
    
    \begin{itemize}
        \item Méthode 1:
        \begin{enumerate}
            \item Si l'entier est positif, on le convertit en binaire sur 8 bits avec la méthode de la question 1.
            \item Si l'entier est négatif,
            \begin{itemize}
                \item On prend la valeur absolue de l'entier.
                \item On le convertit en binaire sur 8 bits avec la méthode de la question 1.
                \item On inverse tous les bits du résultat.
                \item On additionne 1 au résultat précédent.
                \item On obtient le complément à 2 de l'entier.
            \end{itemize}
        \end{enumerate}
        \item Méthode 2:
        \begin{enumerate}
            \item Si l'entier est positif, on le convertit en binaire sur 8 bits avec la méthode de la question 1.
            \item Si l'entier est négatif,
            \begin{itemize}
                \item On prend la valeur absolue de l'entier.
                \item On le convertit en binaire sur 8 bits avec la méthode de la question 1.
                \item On recopie les bits de la droite vers la gauche jusqu'au premier '1' inclus
                \item On inverse tous les bits à gauche du premier '1' inclus.
                \item On obtient le complément à 2 de l'entier.
            \end{itemize}
        \end{enumerate}
    \end{itemize}
    \end{answer}

    \item Par la méthode que vous avez explicité précédemment (question 2.), convertir les valeurs suivantes en binaire sur 8 bits:
    
    \begin{answer}{1}{1.5}
    
        \begin{minipage}[t]{0.33\linewidth}
            \centering{$-24$}

            Réponse: $11101000$
        \end{minipage}
        \begin{minipage}[t]{0.33\linewidth}
            \centering{$-124$}

            Réponse: $10000100$
        \end{minipage}
        \begin{minipage}[t]{0.33\linewidth}
            \centering{$9$}

            Réponse: $00001001$
        \end{minipage}
    
    \end{answer}

\end{enumerate}

\newpage

\thispagestyle{DS_LP}
\exods{\ptsexoC}

Ecrire la fonction \texttt{conversion\_relatif\_binaire} qui convertit un entier relatif en binaire sur 8 bits. Vous utiliserez les fonctons précédentes pour créer votre fonction. 
    
    \begin{answer}{1}{7}
    \begin{verbatim}
def conversion_relatif_binaire(val_entiere:int)->str:
    ''' Convertit un entier relatif en binaire sur 8 bits. 
    val_entiere : int : entier relatif compris entre -128 et 127
    return : str : représentation binaire sur 8 bits de l'entier relatif'''

    if val_entiere >= 0:
        return conversion_entier_binaire(val_entiere)
    else:
        val_entiere =  - val_entiere
        binaire = conversion_entier_binaire(val_entiere)

        binaire_inverse = ''
        for bit in binaire:
            binaire_inverse += inverse_bit(bit)
        return additionne_binaire(binaire_inverse, '00000001')

\end{verbatim}
    \end{answer}

    \begin{answer}{1}{11}
        \begin{verbatim}
    def conversion_relatif_binaire(val_entiere:int)->str:
        ''' Convertit un entier relatif en binaire sur 8 bits. 
        val_entiere : int : entier relatif compris entre -128 et 127
        return : str : représentation binaire sur 8 bits de l'entier relatif'''
    
        if val_entiere >= 0:
            return conversion_entier_binaire(val_entiere)
        else:
            val_entiere =  - val_entiere
            binaire = conversion_entier_binaire(val_entiere)
            
            resultat = ''
            i = len(binaire) - 1
            premier_1_trouve = False
            # On parcourt le binaire de droite à gauche
            while i >= 0:
                bit = binaire[i]

                if not premier_1_trouve:
                    resultat = bit + resultat
                    if bit == '1':
                        premier_1_trouve = True
                    else:
                        resultat = inverse_bit(bit) + resultat
                i -= 1
                  \end{verbatim}
        \end{answer}
    
    
\exods{\ptsexoD}


Entourer la ou les bonne(s) réponse(s)

\begin{enumerate}
    \item Que peut-on dire de l'entier relatif $1010001$ codé sur 8 bits en complément à deux ?
    \begin{enumerate}[A)]
        \begin{minipage}[t]{0.24\linewidth}
        \item Cet entier est pair 
        \end{minipage}\hfill
        \begin{minipage}[t]{0.24\linewidth}
            \item \fbox{Cet entier est négatif}
        \end{minipage}\hfill
        \begin{minipage}[t]{0.24\linewidth}
            \item Cet entier est positif
        \end{minipage}\hfill
        \begin{minipage}[t]{0.24\linewidth}
            \item Cet entier vaut $-12_{10}$
        \end{minipage}
    \end{enumerate}

    \medskip

    \item Si $1011$ est un entier signé sur 4 bits en complément à 2, alors sa valeur décimale est:
        \begin{enumerate}[A)]
        \begin{minipage}[t]{0.24\linewidth}
        \item $-7$
        \end{minipage}\hfill
        \begin{minipage}[t]{0.24\linewidth}
        \item \fbox{$-5$}
        \end{minipage}\hfill
        \begin{minipage}[t]{0.24\linewidth}
        \item $5$       
        \end{minipage}\hfill
        \begin{minipage}[t]{0.24\linewidth}
        \item $7$
        \end{minipage}
    \end{enumerate}
\end{enumerate}



\end{document}








